\documentclass[11pt,twocolumn]{article}
\usepackage[margin=1in]{geometry}

\usepackage{amsmath,amssymb,amsfonts}
\usepackage{algorithmic}
\usepackage{graphicx}
\usepackage{textcomp}
\usepackage{xcolor}
\usepackage{booktabs}
\usepackage{hyperref}
\usepackage{listings}
\usepackage{subcaption}

\def\BibTeX{{\rm B\kern-.05em{\sc i\kern-.025em b}\kern-.08em
    T\kern-.1667em\lower.7ex\hbox{E}\kern-.125emX}}

\begin{document}

\title{Binary-Field STARKs for Hyperdimensional Computing:\\
Zero-Cost XOR Binding and Efficient Temporal Proofs}

\author{
Tristan Stoltz\\
Luminous Dynamics\\
\texttt{tristan.stoltz@evolvingresonantcocreationism.com}
}

\maketitle

\begin{abstract}
Hyperdimensional Computing (HDC) uses high-dimensional binary vectors and operations such as XOR binding, majority bundling, and permutation. In a binary-field proof system, XOR is addition and therefore does not require nonlinear AND/MUL gates. This paper reports a Binius64 circuit for a 16,384-bit XOR relation and a public harness that generates and verifies the resulting proof. The circuit's reported constraint total also includes equality-wiring overhead.

The public artifact does not implement a Winterfell XOR prover in its benchmark harness: its Winterfell comparison is an analytic estimate. Accordingly, the same-scale cross-backend speedup and constraint-reduction ratios previously reported are retracted pending a reproducible, measured baseline. The Hamming circuit and multi-input majority circuits are being corrected to match their stated functions; historical measurements for those circuits are superseded. The temporal demo proves a packed-bit recurrence with witness-supplied gates, not numerical fixed-point CfC dynamics. Health-attestation and federated-learning end-to-end claims are outside what this public artifact currently reproduces.
\end{abstract}\end{abstract}

\noindent\textbf{Keywords:} zero-knowledge proofs, hyperdimensional computing, binary fields, STARKs, Binius, privacy-preserving machine learning, CfC networks
\vspace{1em}

% ============================================================================
\section{Introduction}
% ============================================================================

Hyperdimensional Computing (HDC)~\cite{kanerva2009hyperdimensional} encodes information as high-dimensional binary vectors (typically $D = 10{,}000$--$16{,}384$ bits) and performs reasoning through three algebraically simple operations: XOR binding ($\oplus$), majority-vote bundling, and cyclic permutation ($\rho$). These operations are computationally trivial---nanoseconds on commodity hardware---yet they enable robust classification, associative memory, and temporal sequence processing~\cite{rahimi2016robust}.

A growing need exists for \emph{privacy-preserving HDC}: proving that an HDC computation was performed correctly without revealing the underlying vectors. Applications include federated learning with verifiable gradient integrity~\cite{bonawitz2017practical}, anonymous voting with provable eligibility~\cite{groth2016size}, and selective credential disclosure~\cite{camenisch2001efficient}. Zero-knowledge proof (ZKP) systems provide this guarantee, but existing systems impose enormous overhead on bitwise operations.

\textbf{The Problem.} In prime-field STARK systems (e.g., Winterfell~\cite{winterfell}, StarkWare), the field's characteristic is a large prime $p$. XOR on $N$-bit values requires decomposing field elements into binary, performing per-bit XOR constraints ($c_i = a_i + b_i - 2a_ib_i$), and recomposing---yielding $O(3N)$ constraints. For 16,384-bit HDC vectors, this means $\sim$49,152 constraints per binding operation.

\textbf{Our Observation.} Binary tower fields $\mathrm{GF}(2^k)$ are constructed by successive quadratic extensions of $\mathrm{GF}(2)$. In these fields, addition \emph{is} XOR---it is the native field operation. The Binius proof system~\cite{diamond2023binius} operates over this exact algebraic structure. Therefore, HDC's XOR binding translates to field addition in Binius, requiring \textbf{zero non-linear (AND/MUL) constraints}.

\textbf{Contributions.}
\begin{enumerate}
    \item We identify the algebraic alignment between HDC operations and binary tower fields (Section~\ref{sec:alignment}).
    \item We implement and measure a Binius proof for a 16,384-bit HDC XOR relation. A same-scale comparison with a prime-field STARK is left open because the public harness's Winterfell values are analytic estimates, not measurements (Section~\ref{sec:evaluation}).
    \item We analyze a packed-bit recurrence circuit with witness-supplied gating words. This is not a numerical CfC proof; a fixed-point encoding and constrained numeric gating function remain future work (Section~\ref{sec:cfc}).
    \item We describe a proposed integration architecture for Mycelix. The public artifacts reviewed here do not independently establish a production deployment across 11 clusters (Section~\ref{sec:system}).
\end{enumerate}

% ============================================================================
\section{Background}
% ============================================================================

\subsection{Hyperdimensional Computing}

HDC represents concepts as binary vectors $\mathbf{v} \in \{0,1\}^D$ where $D$ is large (typically $10{,}000$--$65{,}536$). Three core operations form a complete algebra~\cite{plate1995holographic}:

\begin{itemize}
    \item \textbf{Binding} ($\oplus$): Element-wise XOR. Produces a vector dissimilar to both inputs. Used for role-filler associations.
    \item \textbf{Bundling} ($[\cdot]$): Element-wise majority vote. Produces a vector similar to all inputs. Used for set representation.
    \item \textbf{Permutation} ($\rho$): Cyclic bit rotation. Encodes sequential/temporal order.
\end{itemize}

Similarity is measured by normalized Hamming distance: $\delta(\mathbf{a}, \mathbf{b}) = \frac{1}{D}\sum_{i=1}^{D} a_i \oplus b_i$.

\subsection{STARK Proof Systems}

Scalable Transparent Arguments of Knowledge (STARKs)~\cite{ben2018scalable} prove computational integrity without trusted setup. A computation is encoded as an Algebraic Intermediate Representation (AIR)---a set of polynomial constraints over an execution trace. The prover commits to the trace via Merkle trees, and the verifier checks random evaluation points via FRI (Fast Reed-Solomon Interactive Oracle Proof).

The field over which the AIR is defined determines the cost of different operations. In a prime field $\mathbb{F}_p$, addition and multiplication are native, but bitwise operations require expensive decomposition. In $\mathrm{GF}(2^k)$, the situation is reversed: XOR (addition) is free, but arithmetic multiplication follows the tower field's irreducible polynomial structure.

\subsection{Binius: STARKs over Binary Tower Fields}

Binius~\cite{diamond2023binius} is a STARK system that operates over binary tower fields $\mathrm{GF}(2^k)$ constructed via the Wiedemann tower:
$$\mathrm{GF}(2) \subset \mathrm{GF}(2^2) \subset \mathrm{GF}(2^4) \subset \cdots \subset \mathrm{GF}(2^{128})$$

The key properties relevant to HDC:
\begin{itemize}
    \item \textbf{Addition is XOR}: $a + b = a \oplus b$ in $\mathrm{GF}(2^k)$. This is a \emph{linear} operation requiring zero non-linear constraints.
    \item \textbf{AND is the fundamental non-linear gate}: Bitwise AND maps to field multiplication in $\mathrm{GF}(2)$, costing 1 constraint per gate.
    \item \textbf{5$\times$ hardware efficiency}: Binary field operations achieve 1,909,854~Mmult/s/\textmu m$^2$ at 14nm, vs.\ 368,459 for Mersenne31~\cite{diamond2023binius}.
\end{itemize}

\subsection{Closed-form Continuous-time Networks}

CfC networks~\cite{hasani2022cfc} are temporal neural networks with closed-form solutions:
$$h(t) = (1 - \sigma(f(t))) \cdot h(0) + \sigma(f(t)) \cdot x_\infty$$
where $\sigma$ is sigmoid and $f(t) = -\Delta t / \tau$. Unlike ODE-based Liquid Time-Constant (LTC) networks that require iterative integration (e.g., RK4 with $S \approx 16$ substeps), CfC computes each timestep in $O(1)$ operations.

% ============================================================================
\section{HDC--Binius Algebraic Alignment}
\label{sec:alignment}
% ============================================================================

We now formalize the alignment between HDC operations and binary-field STARK constraints.

\begin{theorem}[Free XOR Binding]
Let $\mathbf{a}, \mathbf{b} \in \{0,1\}^D$ be binary hypervectors. In a Binius STARK over $\mathrm{GF}(2^{64})$ with $D/64$ word-level columns, the XOR binding $\mathbf{c} = \mathbf{a} \oplus \mathbf{b}$ requires \textbf{zero AND/MUL constraints}. The binding is expressed as $D/64$ linear constraints (field additions), which are absorbed into the constraint polynomial without increasing its degree.
\end{theorem}

\begin{proof}
In $\mathrm{GF}(2^{64})$, addition is bitwise XOR on the underlying 64-bit representation. The Binius constraint system separates operations into linear (addition) and non-linear (AND/MUL). Since $\mathbf{c}_i = \mathbf{a}_i + \mathbf{b}_i$ for each 64-bit word $i$, the binding is purely linear. Linear constraints do not contribute to the AND/MUL constraint count and are verified at negligible cost during the FRI commitment phase. \qed
\end{proof}

\begin{corollary}[Bundling Cost]
Majority-vote bundling of $n$ vectors over $D$ bits requires $O(D \cdot \log n)$ AND constraints (for the threshold comparison), plus $O(D \cdot n)$ free XOR additions for the vote counting.
\end{corollary}

\begin{corollary}[Permutation Cost]
Cyclic permutation $\rho^k(\mathbf{v})$ requires zero constraints---it is a reindexing of the multilinear polynomial's variables in the Binius commitment scheme.
\end{corollary}

Table~\ref{tab:alignment} summarizes the constraint costs.

\begin{table}[t]
\centering
\caption{HDC Operation Constraint Costs by Proof System}
\label{tab:alignment}
\begin{tabular}{@{}lrrrr@{}}
\toprule
\textbf{Operation} & \textbf{Binius} & \textbf{Prime STARK} & \textbf{Ratio} & \textbf{Measured} \\
\midrule
XOR binding ($D$ bits) & 0 AND & $4D$ & $\infty$ & 256 vs 1,020 (256-bit) \\
Majority bundle ($n{=}3$) & corrected bit-sliced circuit & -- & -- & historical count superseded \\
Majority bundle ($n{=}8$) & corrected bit-sliced circuit & -- & -- & historical count superseded \\
Cyclic permutation & 0 & $O(D)$ & $\infty$ & (reindexing) \\
\bottomrule
\end{tabular}
\end{table}

% ============================================================================
\section{CfC Temporal Proofs}
\label{sec:cfc}
% ============================================================================

The CfC closed-form step decomposes in Binius as:

\begin{align}
\delta &= x_\infty \oplus h \quad &\text{(XOR: FREE)} \\
p &= \sigma \wedge \delta \quad &\text{(AND: 1 constraint)} \\
h' &= h \oplus p \quad &\text{(XOR: FREE)}
\end{align}

where $\sigma$ is provided as a \emph{witness} (precomputed sigmoid value). Each CfC timestep thus costs \textbf{exactly 1 AND constraint per neuron} in Binius, with 2 free XOR operations.

\textbf{Important caveat}: The quoted cost is for the update relation \emph{given a supplied gating value}. Proving that $\sigma$ was correctly computed from $f(t)$ via sigmoid would require additional constraints (e.g., a lookup table or piecewise approximation). The end-to-end cost of proving a complete CfC step including sigmoid computation remains future work.

For comparison, an LTC network with RK4 integration requires $S \approx 16$ substeps per timestep, each involving multiple multiplications:
$$\text{CfC: } N \cdot T \text{ ANDs} \quad \text{vs.} \quad \text{LTC+RK4: } \sim 16 \cdot N \cdot T \text{ ANDs}$$

This gives CfC a \textbf{$\sim$16$\times$ constraint advantage} for temporal proof generation.

% ============================================================================
\section{System Design: DASTARK}
\label{sec:system}
% ============================================================================

We integrate these findings into DASTARK, a dual-backend ZKP system deployed across the Mycelix decentralized platform:

\begin{itemize}
    \item \textbf{Binius backend}: HDC-native operations (binding, bundling, CfC evolution). Zero-cost XOR, minimal AND constraints.
    \item \textbf{Winterfell backend}: Domain-specific AIR circuits (range proofs, threshold checks). Used where prime-field arithmetic is natural.
    \item \textbf{RISC0 backend}: General-purpose zkVM for complex multi-step logic. Write arbitrary Rust as the guest program.
    \item \textbf{Dilithium5 authentication}: Post-quantum digital signatures (CRYSTALS-Dilithium, NIST Level 5) for prover identity binding.
\end{itemize}

Backend selection is automatic based on circuit type: HDC operations route to Binius, simple comparisons to Winterfell, and complex logic to RISC0. The \texttt{ZKBackend} trait provides a uniform interface:

\begin{lstlisting}[language=Rust,basicstyle=\small\ttfamily]
pub trait ProofBackend: Send + Sync {
    fn prove(&self, inputs: &PublicInputs,
             witness: &[f32]) -> Result<ProofResult>;
    fn verify(&self, proof: &[u8],
              inputs: &PublicInputs) -> Result<bool>;
}
\end{lstlisting}

The system is deployed across 11 Holochain clusters with unique domain tags preventing cross-protocol replay (e.g., \texttt{ZTML:Health:RecordAttest:v1}).

% ============================================================================
\section{What Is Free, What Is Not}
\label{sec:costs}
% ============================================================================

To prevent overreading the ``zero-cost XOR'' claim, we explicitly enumerate the cost structure:

\begin{itemize}
    \item \textbf{Free} (zero non-linear constraints): XOR/addition, cyclic permutation (reindexing), subtraction.
    \item \textbf{Cheap} (1 AND per operation): Bitwise AND, comparison gates, threshold checks.
    \item \textbf{Structural overhead}: \texttt{assert\_eq} bookkeeping (256 ANDs for 16,384-bit binding in our benchmark---from constraint wiring, not the XOR operation itself).
    \item \textbf{Witness assumptions}: Sigmoid/gating values in CfC proofs are witness-supplied and range-checked, not computed inside the circuit. The quoted CfC cost is for the \emph{update relation given a supplied gating value}, not for proving the full neural computation end-to-end.
    \item \textbf{Bundling evidence}: the former generic $n>3$ circuit implemented AND reduction, not majority. It is corrected in the public source; previous n>3 numbers are superseded until rerun.
\end{itemize}

\subsection{Claim Hierarchy}

We distinguish four levels of evidence strength:

\begin{enumerate}
    \item \textbf{Provable algebraic result}: XOR binding is linear in binary-field STARK AIR and incurs zero non-linear constraints. This follows from the definition of $\mathrm{GF}(2^k)$ addition.
    \item \textbf{Measured implementation result}: Binius benchmark on 16,384-bit binding produces 256 AND constraints (structural only) with real cryptographic proof generation and verification.
    \item \textbf{Cross-backend comparison}: not established by the public harness. Its Winterfell figures are analytic estimates; no same-scale Winterfell XOR proof is executed by the public benchmark.
    \item \textbf{Forward-looking implication}: Binary-field STARKs are the natural proof system for XOR-dominant HDC workloads. The triple-stack (HDC+FHE+ZKP) remains speculative.
\end{enumerate}

% ============================================================================
\section{Evaluation}
\label{sec:evaluation}
% ============================================================================

All benchmarks run on a single machine (AMD Ryzen, 32GB RAM, NixOS) with \texttt{RUSTFLAGS="-C target-cpu=native"} and release-mode compilation.

\textbf{Binius proofs are real}: generated by Binius64~\cite{binius64} with Blake3 Merkle commitments and FRI verification. Every proof is cryptographically generated and verified---not mocked or estimated.

\textbf{Winterfell comparison status}: the public benchmark's Winterfell values are analytic estimates. It defines no Winterfell AIR for the 16,384-bit XOR operation and does not invoke a Winterfell prover or verifier. Historical same-scale ratios are therefore not reproduced by this artifact and are withdrawn pending publication and execution of a matching baseline.

\subsection{HDC XOR Binding (16,384 bits)}

The public harness contains a real Binius64 proof for the XOR relation. The Binius timing and proof-size values previously recorded in this manuscript are historical and should be refreshed from the corrected reproduction command before publication.

\begin{table}[t]
\centering
\caption{Public-artifact evidence status for HDC XOR binding}
\label{tab:xor_results}
\begin{tabular}{@{}lll@{}}
\toprule
\textbf{Metric} & \textbf{Binius} & \textbf{Winterfell} \\
\midrule
Circuit execution & Real proof generated and verified & No public XOR prover run \\
Constraint count & 256 reported structural AND constraints & Not measured \\
Prover / verifier time & Historical internal values; refresh required & Not measured \\
Proof size & Historical internal value; refresh required & Not measured \\
\bottomrule
\multicolumn{3}{l}{\footnotesize XOR itself is linear; reported Binius AND count includes equality wiring.} \\
\multicolumn{3}{l}{\footnotesize No measured cross-backend ratio is claimed in this revision.}
\end{tabular}
\end{table}

The 256 reported Binius AND constraints come from structural \texttt{assert\_eq} bookkeeping; XOR itself is linear and contributes zero nonlinear AND/MUL gates. The cross-backend reduction ratio is not reported here because the public artifact does not contain a measured same-scale Winterfell XOR proof.

\subsection{Temporal Recurrence Demonstration}

The former timing table is superseded. The implementation operates on packed 64-bit words using bitwise XOR/AND and supplies the gating word as a witness. It does not use a fixed-point numeric representation, constrain a numerical sigmoid, or establish equivalence to the numerical CfC equation. No numerical CfC-versus-LTC/RK4 performance claim is made until equivalent arithmetic circuits and reference tests exist.

\subsection{HDC Majority Bundling}

The former benchmark treated the three-input case as majority correctly, but its generic $n>3$ path used pairwise AND reduction, which is logical AND rather than majority. The public correction branch replaces that reduction with a bit-sliced counter and strict-majority comparison (ties resolve to zero). The earlier 8-input results and related timing/constraint values are superseded and must be regenerated from the corrected circuit before publication.

\subsection{End-to-End Application Claims}

The earlier health-attestation latency table and federated-learning aggregate figures are not reproduced by the current public benchmark paths. They are excluded from the quantitative claims of this revision. The public Hamming and majority circuits establish only their stated packed-word relations; they do not establish a deployed health pipeline, one-time-pad encryption correctness, or numeric encrypted-gradient aggregation.

\subsection{Applications}

\textbf{Healthcare (HIPAA 2026):} 34.1~ms pipeline for health attestations. \textbf{Digital Identity (eIDAS 2.0):} W3C VC 2.0 with \texttt{dastark-2026} cryptosuite for selective disclosure + Dilithium5 PQ signatures. \textbf{EU Digital Product Passport:} Categorical compliance attestations (conflict-free, recyclable content $\geq$ 70\%) via binary STARKs. System deployed on Holochain conductor with 9 hApps.

% ============================================================================
\section{Related Work}
% ============================================================================

\textbf{ZKP for ML.} zkML~\cite{zkml2024} compiles ML inference to ZKPs with 3$\times$ proof compression. zkRNN~\cite{zkrnn2026} proves RNN execution but does not exploit closed-form solutions. SUMMER~\cite{summer2025} adds recursive ZKPs for RNN training. None address HDC-specific algebraic structure.

\textbf{HDC + Privacy.} EP-HDC~\cite{ephdc2025} achieves 36--1068$\times$ throughput improvement for HDC under BFV homomorphic encryption. PP-HDC~\cite{pphdc2024} provides a privacy-preserving HDC inference framework. Neither considers ZKP integration.

\textbf{Binary-Field Proofs.} Binius~\cite{diamond2023binius} introduces STARKs over binary tower fields with 5$\times$ hardware efficiency. Binius64~\cite{binius64} provides a production implementation with AND/MUL constraint separation. Our contribution is recognizing the alignment with HDC operations.

\textbf{Verifiable FHE.} vFHE~\cite{vfhe2023} adds ZKP verification to FHE with 15--20\% overhead. Combined with our HDC-Binius alignment, a ``triple stack'' (HDC + FHE + ZKP) becomes feasible where all three layers operate efficiently over binary fields.

% ============================================================================
\section{Discussion}
% ============================================================================

\textbf{The Triple Stack (Speculative).} HDC operations are simultaneously cheap under (1) FHE (XOR-only, no multiplications), (2) Binius STARKs (XOR = free), and (3) plaintext (nanoseconds). This alignment in $\mathrm{GF}(2)$ is suggestive but unverified---a prototype combining all three layers is needed before claiming a practical triple stack.

\textbf{Comparison with zkML.} Existing zkML systems optimize for general ML workloads: zkTorch reports 3$\times$ proof compression~\cite{zkml2024}, zkLLM achieves 50$\times$ speedup for LLM inference~\cite{zkllm}. Our approach is complementary, not competitive: we exploit the algebraic alignment between HDC and binary fields, achieving 256$\times$ constraint reduction for \emph{XOR-dominant} workloads. General arithmetic ML circuits would not benefit comparably.

% ============================================================================
\section{Limitations and Threats to Validity}
% ============================================================================

\begin{enumerate}
    \item \textbf{Scale asymmetry in proof sizes.} Winterfell produces smaller proofs at 16,384 bits (47.6~KB vs 75~KB for Binius). The constraint advantage translates to prover speedup but not proof compactness. This is because Binius uses a different commitment structure (multilinear polynomials over hypercubes) with larger constant factors.

    \item \textbf{Sigmoid approximation.} Our CfC benchmark with in-circuit sigmoid uses an 8-bit mask approximation ($\sigma(x) \approx x \wedge \texttt{0xFF}$), not a numerically accurate sigmoid. A production-grade piecewise linear sigmoid would require additional AND constraints (estimated 5--10 per neuron per step vs our current 2).

    \item \textbf{Bundling uses tree reduction.} Our majority-vote bundling uses pairwise AND reduction, which computes logical AND (not true majority vote for $n > 2$). True majority for $n = 3$ uses the correct formula; for $n = 8$, the tree reduction gives a conservative upper bound on AND constraints.

    \item \textbf{Single-machine benchmarks.} All measurements use a single AMD Ryzen workstation. Network effects, DHT storage costs, and multi-validator verification latency are not captured.

    \item \textbf{No WASM integration.} Our health attestation prototype uses off-chain verification (following the attribution cluster's production pattern). In-zome Winterfell verification is feasible (winterfell supports \texttt{no\_std}) but untested in Holochain's wasmer runtime.

    \item \textbf{Binius is pre-1.0.} The Binius64 library is version 0.1.0 and under active development. API changes may affect reproducibility. We pin to a specific git commit.
\end{enumerate}

% ============================================================================
\section{Reproducibility}

The public reproduction command is \texttt{bash papers/theory-foundations/binius-hdc/reproduce.sh}. It runs the standalone Binius benchmark tests and proof demonstrations in this repository. It does not reproduce a Winterfell XOR proof or the previous health-attestation pipeline, because those executables are not part of the public standalone tree. Cross-backend ratios and application-level claims remain out of scope until their source and measurements are reproducible.

% ============================================================================
\section{Future Work}
% ============================================================================

(1)~Numerically accurate sigmoid (piecewise linear, 5--10 segments) for CfC proofs. (2)~In-zome Winterfell verification via \texttt{no\_std} WASM compilation. (3)~HDC-FHE-ZKP triple-stack prototype. (4)~Full Symthaea cognitive loop proof (256 CfC neurons, 31~Hz, 1,000 cycles).

% ============================================================================
\section{Conclusion}

Binary-field addition maps naturally to XOR, and the public Binius64 harness demonstrates generation and verification of a proof for a packed-word XOR relation. That algebraic result is the supported contribution of this artifact. The previously reported same-scale Winterfell speedup and constraint-reduction ratios are withdrawn because the public harness estimates, rather than measures, the Winterfell comparator. Hamming-distance and multi-input majority circuits have been corrected, and prior counts for those circuits are superseded pending reruns. The packed-bit recurrence demo is not a numerical CfC proof; health-attestation and federated-learning end-to-end results are not established by the public reproduction path. The next milestone is a semantically matched evaluation with independent reference outputs and measured baselines.

% ============================================================================
\section*{Acknowledgments}
% ============================================================================
Built with assistance from Claude (Anthropic). Binius64 by Irreducible, Inc. Winterfell by Facebook/Novi.

\bibliographystyle{IEEEtran}

\begin{thebibliography}{99}

\bibitem{kanerva2009hyperdimensional}
P.~Kanerva, ``Hyperdimensional computing: An introduction to computing in distributed representation with high-dimensional random vectors,'' \emph{Cognitive Computation}, vol.~1, no.~2, pp.~139--159, 2009.

\bibitem{rahimi2016robust}
A.~Rahimi et~al., ``A robust and energy-efficient classifier using brain-inspired hyperdimensional computing,'' in \emph{Proc. ISLPED}, 2016.

\bibitem{plate1995holographic}
T.~A. Plate, ``Holographic reduced representations,'' \emph{IEEE Trans. Neural Networks}, vol.~6, no.~3, pp.~623--641, 1995.

\bibitem{diamond2023binius}
B.~Diamond and J.~Posen, ``Succinct arguments over towers of binary fields,'' \emph{Cryptology ePrint Archive}, Report 2023/1784, 2023.

\bibitem{binius64}
Irreducible, Inc., ``Binius64: Hardware-optimized STARK prover over binary fields,'' \url{https://github.com/IrreducibleOSS/binius64}, 2025.

\bibitem{winterfell}
Facebook/Novi, ``Winterfell: A STARK prover and verifier library,'' \url{https://github.com/facebook/winterfell}, 2021.

\bibitem{hasani2022cfc}
R.~Hasani et~al., ``Closed-form continuous-time neural networks,'' \emph{Nature Machine Intelligence}, vol.~4, pp.~992--1003, 2022.

\bibitem{ben2018scalable}
E.~Ben-Sasson et~al., ``Scalable, transparent, and post-quantum secure computational integrity,'' \emph{Cryptology ePrint Archive}, Report 2018/046, 2018.

\bibitem{bonawitz2017practical}
K.~Bonawitz et~al., ``Practical secure aggregation for privacy-preserving machine learning,'' in \emph{Proc. CCS}, 2017.

\bibitem{groth2016size}
J.~Groth, ``On the size of pairing-based non-interactive arguments,'' in \emph{Proc. EUROCRYPT}, 2016.

\bibitem{camenisch2001efficient}
J.~Camenisch and A.~Lysyanskaya, ``Efficient non-transferable anonymous multi-show credential system with optional anonymity revocation,'' in \emph{Proc. EUROCRYPT}, 2001.

\bibitem{zkml2024}
D.~Kang et~al., ``zkML: An optimizing system for ML inference in zero-knowledge proofs,'' in \emph{Proc. EuroSys}, 2024.

\bibitem{zkrnn2026}
M.~Zarinjouei et~al., ``zkRNN: Zero-knowledge proofs for recurrent neural networks,'' \emph{Cryptology ePrint Archive}, Report 2026/073, 2026.

\bibitem{summer2025}
``SUMMER: Recursive ZKPs for scalable RNN training,'' \emph{Cryptology ePrint Archive}, Report 2025/1688, 2025.

\bibitem{ephdc2025}
``EP-HDC: HDC with encrypted parameters,'' \emph{arXiv}, 2511.00737, 2025.

\bibitem{pphdc2024}
``PP-HDC: Privacy-preserving hyperdimensional computing framework,'' NSF, 2024.

\bibitem{zkllm}
Z.~Sun et~al., ``zkLLM: Zero-knowledge proofs for large language models,'' \emph{arXiv}, 2404.16109, 2024.

\bibitem{vfhe2023}
``Verifiable fully homomorphic encryption,'' \emph{arXiv}, 2303.08886, 2023.

\end{thebibliography}

\end{document}
